\documentclass[11pt]{article}
\usepackage[a4paper,margin=1in]{geometry}
\usepackage{graphicx}
\usepackage{listings}
\usepackage{xcolor}
\usepackage{booktabs}
\usepackage{array}
\usepackage{hyperref}
\usepackage{parskip}
\hypersetup{colorlinks=true, linkcolor=black, urlcolor=blue}

\definecolor{codebg}{RGB}{245,245,245}
\lstdefinestyle{cstyle}{
  language=C,
  backgroundcolor=\color{codebg},
  basicstyle=\ttfamily\footnotesize,
  keywordstyle=\color{blue}\bfseries,
  commentstyle=\color{gray}\itshape,
  stringstyle=\color{teal},
  numbers=left,
  numberstyle=\tiny\color{gray},
  stepnumber=1,
  breaklines=true,
  frame=single,
  tabsize=2,
  showstringspaces=false
}
\lstdefinestyle{outstyle}{
  backgroundcolor=\color{codebg},
  basicstyle=\ttfamily\footnotesize,
  breaklines=true,
  frame=single,
  showstringspaces=false
}
\lstdefinestyle{bashstyle}{
  language=bash,
  backgroundcolor=\color{codebg},
  basicstyle=\ttfamily\footnotesize,
  breaklines=true,
  frame=single,
  showstringspaces=false
}

\title{\vspace{-2cm}}
\date{}

\begin{document}
\begin{titlepage}
\centering
\vspace*{1cm}
{\large \textbf{Sri Sivasubramaniya Nadar College of Engineering}}\\
{\small (An Autonomous Institution, Affiliated to Anna University, Chennai)}\\
{\small Kalavakkam -- 603 110}\\[0.3cm]
{\small Department of Computer Science \& Engineering}\\[1.5cm]

\rule{\linewidth}{0.5mm}\\[0.4cm]
{\LARGE \textbf{Computer Networks Laboratory}}\\[0.3cm]
{\Large 1513 -- Computer Networks Lab}\\[0.4cm]
\rule{\linewidth}{0.5mm}\\[1.5cm]

{\huge \textbf{Assignment 3}}\\[0.4cm]
{\Large Implementation of HTTP Client using TCP Sockets}\\[2cm]

\begin{center}
\begin{tabular}{ll}
\textbf{Name:} & Sundareswaran Rajaguru \\[0.2cm]
\textbf{Register Number:} & 3122247001048 \\[0.2cm]
\textbf{Year / Semester:} & III Year / V Semester \\[0.2cm]
\textbf{Course:} & 1513 Computer Networks Lab \\[0.2cm]
\textbf{Date of Submission:} & 22 July 2026 \\[0.2cm]
\end{tabular}
\end{center}

\vfill
\end{titlepage}

\section*{1. Objective}
To develop a TCP socket-based HTTP client in C that accepts a URL from the user, parses the
host/port/path, establishes a TCP connection with a web server, sends an HTTP GET request,
retrieves the requested resource, and stores it locally. The client communicates with a
Python-based HTTP server (\texttt{http.server}) hosting HTML, text, PDF and image files, and
analyses the communication performance in terms of response time, download time, data size and
throughput.

\section*{2. Server Algorithm (Python HTTP Server)}
The assignment specification requires the server side to be a stock Python web server rather
than custom C code, so that the C client can be validated against a standards-compliant HTTP
server.
\begin{enumerate}
  \item Create a directory \texttt{CN\_Web\_Server} and place the resources to be served inside
        it: \texttt{index.html}, \texttt{notes.txt}, \texttt{sample.pdf}, \texttt{image.jpg}.
  \item Start Python's built-in HTTP server bound to TCP port 8000:
  \begin{lstlisting}[style=bashstyle]
python3 -m http.server 8000
  \end{lstlisting}
  \item Internally, \texttt{http.server} creates a TCP listening socket, \texttt{accept()}s each
        client connection, parses the HTTP request line, maps the requested path to a file in the
        served directory, and streams back a well-formed \texttt{HTTP/1.0 200 OK} response with
        \texttt{Content-Type} and \texttt{Content-Length} headers followed by the file bytes (or a
        \texttt{404} if the resource is absent).
  \item Every request/response pair is logged to standard output (source address, timestamp,
        request line, status code).
\end{enumerate}

\section*{3. Client Algorithm}
\begin{enumerate}
  \item Prompt the user for a URL and parse it into \texttt{host}, \texttt{port} (default 80) and
        \texttt{path} (\texttt{parse\_url()}).
  \item Resolve the host name using \texttt{gethostbyname()}.
  \item Create a TCP socket and \texttt{connect()} to \texttt{host:port}.
  \item Build and \texttt{send()} an HTTP/1.1 \texttt{GET} request with a \texttt{Connection: close}
        header so the server closes the connection once the response is complete.
  \item Record a start timestamp, then loop on \texttt{recv()}:
  \begin{itemize}
    \item Record the timestamp of the first received byte (used for \emph{response time}).
    \item Parse and discard the HTTP response header (up to the blank line \texttt{\textbackslash r\textbackslash n\textbackslash r\textbackslash n}).
    \item Write the remaining body bytes to a local output file named
          \texttt{downloaded\_<resource>}.
  \end{itemize}
  \item Record an end timestamp once \texttt{recv()} returns 0 (server closed the connection).
  \item Compute and display performance metrics: response time, download time, data size and
        throughput (KB/s).
\end{enumerate}

\section*{4. Server Code}
The server is Python's standard-library \texttt{http.server} module used as-is (no custom source
required by the assignment); it is launched with the single command shown in Section~2.

\section*{5. Client Code}
\subsection*{5.1 http\_client.c}
\lstinputlisting[style=cstyle]{http_client.c}

\section*{6. Output / Test Cases}
The Python server was started inside \texttt{CN\_Web\_Server} on port 8000 and the C client was
run four times, once for each resource type required by the assignment: an HTML page, a plain
text file, a PDF document and a JPEG image.

\subsection*{6.1 Server-side Access Log}
\begin{lstlisting}[style=outstyle]
::1 - - [17/Aug/2026 10:42:49] "HEAD /index.html HTTP/1.1" 200 -
::1 - - [17/Aug/2026 10:43:48] "GET /index.html HTTP/1.1" 200 -
::ffff:127.0.0.1 - - [17/Aug/2026 10:43:55] "GET /index.html HTTP/1.1" 200 -
::ffff:127.0.0.1 - - [17/Aug/2026 10:43:55] "GET /notes.txt HTTP/1.1" 200 -
::ffff:127.0.0.1 - - [17/Aug/2026 10:43:55] "GET /sample.pdf HTTP/1.1" 200 -
::ffff:127.0.0.1 - - [17/Aug/2026 10:43:55] "GET /image.jpg HTTP/1.1" 200 -
\end{lstlisting}

\subsection*{6.2 Client Console Output (all four resource types)}
\lstinputlisting[style=outstyle]{output.txt}

\subsection*{6.3 Performance Summary}
\begin{center}
\renewcommand{\arraystretch}{1.2}
\begin{tabular}{lrrrr}
\toprule
\textbf{Resource} & \textbf{Response (ms)} & \textbf{Download (ms)} & \textbf{Size (bytes)} & \textbf{Throughput (KB/s)} \\
\midrule
index.html & 0.935 & 0.151 & 168   & 1086.51 \\
notes.txt  & 0.568 & 0.447 & 159   & 347.37 \\
sample.pdf & 0.492 & 0.668 & 33790 & 49398.27 \\
image.jpg  & 0.383 & 0.220 & 332   & 1473.72 \\
\bottomrule
\end{tabular}
\end{center}
All four requests succeeded with HTTP \texttt{200 OK}, and each resource was correctly saved
locally as \texttt{downloaded\_<name>}, matching the size reported by the server. The PDF
(largest file) shows the highest throughput because its transfer amortises the fixed
per-connection overhead (TCP handshake, header parsing) over far more payload bytes than the
small HTML/text/image files.

\section*{7. Learning Outcomes}
\begin{itemize}
  \item Learned to manually construct and parse HTTP/1.1 request and response messages at the
        socket level, including correctly detecting the header/body boundary
        (\texttt{\textbackslash r\textbackslash n\textbackslash r\textbackslash n}).
  \item Understood the role of the \texttt{Connection: close} header in enabling a client to
        detect end-of-response via \texttt{recv()} returning 0, without a \texttt{Content-Length}.
  \item Practised interoperating a custom C TCP client with a standards-compliant third-party HTTP
        server (Python's \texttt{http.server}), reinforcing that HTTP is a text-based
        application-layer protocol independent of implementation language.
  \item Learned to measure network application performance metrics -- response time, download
        time, and throughput -- using \texttt{gettimeofday()} at precise points in the data path.
  \item Observed how payload size affects effective throughput due to fixed per-request overhead.
\end{itemize}

\end{document}
